% BHU PYQ -- Manually transcribed & error-corrected
% Paper: BCH-223 : Specialised Accounts-II
% Exam: B.Com. (Hons.) (Semester IV) (Special) Examination, 2014-15

\documentclass[11pt,a4paper]{article}
\usepackage[a4paper,top=2.5cm,bottom=2.5cm,left=2.8cm,right=2.8cm,headheight=14pt]{geometry}
\usepackage{amsmath,amssymb}
\usepackage[T1]{fontenc}
\usepackage[utf8]{inputenc}
\usepackage{lmodern,microtype,fancyhdr,enumitem,hyperref,lastpage,parskip}

\hypersetup{
    colorlinks=true,
    urlcolor=black,
    linkcolor=black,
    pdftitle={BCH-223: Specialised Accounts-II -- BHU B.Com. (Hons.) Sem IV 2014-15}
}

\pagestyle{fancy}
\fancyhf{}
\renewcommand{\headrulewidth}{0.4pt}
\renewcommand{\footrulewidth}{0.4pt}
\fancyhead[L]{\small   \textit{Banaras Hindu University}}
\fancyhead[R]{\small   \textit{BCH-223: Specialised Accounts-II}}
\fancyfoot[C]{\small Page  \thepage\ of \pageref{LastPage}}


\newlist{parts}{enumerate}{2}
\setlist[parts,1]{label=(\roman*),leftmargin=*,itemsep=5pt,topsep=3pt}
\setlist[parts,2]{label=(\alph*),leftmargin=*,itemsep=3pt,topsep=2pt}

\newcommand{\pts}[1]{\hfill   \textit{[#1]}}


\begin{document}


\begin{center}
  {\large\bfseries BANARAS HINDU UNIVERSITY}\\[2pt]
  {
\normalsize B.Com. (Hons.) --- Semester IV (Special) Examination, 2014-15}\\[6pt]
  \rule{0.8\linewidth}{0.5pt}\\[6pt]
  {\large\bfseries Paper No. BCH-223: Specialised Accounts-II}\\[4pt]
  {
\normalsize Time: 3 Hours\hfill Maximum Marks: 70}
\end{center}

\rule{\linewidth}{0.4pt}

\smallskip



\noindent   \textit{Instructions: Answer all questions. All questions carry equal marks (14 marks each).}

\bigskip




\noindent   \textbf{Question 1.}
Prepare the Profit and Loss Account of the   \textit{Mahamana Bank Ltd.} for the year ending 31st March, 2014 from the following particulars:


\begin{center}
\small

\begin{tabular}{|l|r|}
\hline
   \textbf{Particulars} &   \textbf{Amount (Rs.)} \\
\hline
Interest on Loans & 2,59,000 \\
Interest on Fixed Deposits & 2,75,000 \\
Rebate on Bills Discounted & 49,000 \\
Commission Charged to Customers & 8,200 \\
Salaries, etc. & 54,000 \\
Discount on Bills Discounted & 1,95,000 \\
Interest on Cash Credit Accounts & 2,23,000 \\
Interest on Current Accounts & 42,000 \\
Rent and Taxes & 18,000 \\
Interest on Overdrafts & 54,000 \\
Auditors' Fees & 4,200 \\
Interest on Saving Bank Deposits & 68,000 \\
Postages and Telegrams & 1,400 \\
Printing and Stationery & 2,900 \\
Sundry Charges & 1,700 \\
\hline
\end{tabular}
\end{center}\pts{14}

\centerline{   \textbf{OR}}

Give a specimen of a Bank's Profit and Loss Account and explain each head in brief.\pts{14}

\medskip\rule{\linewidth}{0.2pt}

\pagebreak




\noindent   \textbf{Question 2.}
Prepare the Revenue Account in respect of   \textit{Kashi Accident Insurance} for the year ending 31st March, 2014 from the books of General Insurance Ltd.:


\begin{center}
\small

\begin{tabular}{|l|r|}
\hline
   \textbf{Particulars} &   \textbf{Amount (Rs.)} \\
\hline
Reserve for Unexpired Risk (Opening) & 80,000 \\
Additional Reserve & 10,000 \\
Premium (Net) & 5,20,000 \\
Claims (Net) & 1,20,000 \\
Expenses of Management & 10,000 \\
Bonus in Reduction of Premium & 5,000 \\
Commission & 15,000 \\
Outstanding Claims (Closing) & 10,000 \\
Interest and Dividends & 15,000 \\
Income Tax thereon & 4,000 \\
Interest on Investments & 5,000 \\
\hline
\end{tabular}
\end{center}
Provide for Reserve for Unexpired Risks at 50\% of Net premium.\pts{14}

\centerline{   \textbf{OR}}

Give a specimen of the Revenue Account and Balance Sheet of a Life Insurance Company.\pts{14}

\medskip\rule{\linewidth}{0.2pt}




\noindent   \textbf{Question 3.}
Differentiate between the Presidency Towns Insolvency Act and the Provincial Insolvency Act, and also explain the preferential creditors under both Acts.\pts{14}

\centerline{   \textbf{OR}}

Give the proforma of a Statement of Affairs and Deficiency Account using imaginary figures.\pts{14}

\medskip\rule{\linewidth}{0.2pt}

\pagebreak




\noindent   \textbf{Question 4.}
What is Extension and Replacement of Fixed Assets? Explain the accounting treatments regarding extension and replacement of fixed assets.\pts{14}

\centerline{   \textbf{OR}}

The figures given below relate to the   \textit{Maruti Colliery Company} for the year ending 31st March, 2014. You are required to construct a Capital Account and General Balance Sheet under the Double Account System:


\begin{center}
\footnotesize

\begin{tabular}{|l|r||l|r|}
\hline
   \textbf{Particulars} &   \textbf{Amount (Rs.)} &   \textbf{Particulars} &   \textbf{Amount (Rs.)} \\
\hline
Equity Shares (Rs. 10 each) & 2,20,000 & Cash in Hand \& Bank & 60,000 \\
Preference Shares (Rs. 100 each) & 1,20,000 & Office Building & 8,000 \\
Debentures & 60,000 & Furniture & 18,000 \\
Bills Payable & 12,000 & Depreciation Fund & 40,000 \\
Sundry Creditors & 18,000 & Reserve Fund & 30,000 \\
Land & 17,000 & Balance of P \& L A/c (Cr.) & 50,000 \\
Shaft Sinking, etc. & 2,20,000 & Stock & 24,000 \\
Plant \& Machinery & 70,000 & Investment & 34,000 \\
Wagons & 27,000 & Sundry Debtors & 70,000 \\
& & Bills Receivable & 2,000 \\
\hline
\end{tabular}
\end{center}

The above figures include an issue of Rs. 20,000 Preference shares during the year, and the following capital expenditures were made during the same period: Shaft sinking Rs. 12,000; Plant \& Machinery Rs. 10,000; Wagons Rs. 6,000; and Furniture Rs. 2,000.\pts{14}

\medskip\rule{\linewidth}{0.2pt}




\noindent   \textbf{Question 5.}
What do you mean by Government Accounting? Explain the fundamental principles of Government Accounting.\pts{14}

\centerline{   \textbf{OR}}

Differentiate between Commercial Accounting and Government Accounting. Discuss the role of the Comptroller and Auditor General (C\&AG) of India regarding Government Accounting.\pts{14}

\vfill
\rule{\linewidth}{0.4pt}

\begin{center}\small
\href{https://bhu-exam.vercel.app/}{Click here to visit our website}\\[4pt]
\href{https://me-aryan.vercel.app/}{Click here to visit our contributor}\\[4pt]
\href{https://quashan-me.vercel.app/}{Click here to visit our contributor}
\end{center}

\end{document}
